\documentclass[12pt]{article}
\usepackage[utf8]{inputenc}
\usepackage[T1]{fontenc}
\usepackage{hyperref}
\title{Highland Popoluca Language Resources for K-12 Curriculum Development}
\author{Compiled by Claude Sonnet 5.5 (claude-sonnet-5-5) for Andr\'as Kornai}
\date{}
\begin{document}
\maketitle

\section{Introduction}
Highland Popoluca (ISO 639-3 \texttt{poi}; also Sierra Popoluca, Soteapanec) is a Zoquean language of Veracruz, Mexico, with about 36,000 speakers in the 2020 census according to Wikipedia. Existing materials are mostly linguistic, produced by SIL linguists.

\section{Primary Educational Resources}
No primers or school textbooks were found in search results.

\section{Linguistic and Reference Materials}
\begin{itemize}
\item Elson and Guti\'errez, \emph{Diccionario Popoluca de la Sierra, Veracruz} (1999, SIL, 190 pp.), bibliographic listing. \url{https://endangeredlanguages.com/lang/6169/bibliography}
\item Fern\'andez de Miranda and Elson, \emph{Gram\'atica del popoluca de la Sierra} (1960), listed in the same bibliography. \url{https://endangeredlanguages.com/lang/6169/bibliography}
\item B.~F. Elson, ``Sierra Popoluca'', Handbook of Middle American Indians vol.~5 (1967), Glottolog reference. \url{https://glottolog.org/resource/reference/id/16610}
\item Glottolog entry for Highland Popoluca. \url{https://glottolog.org/resource/languoid/id/high1276}
\item Wikipedia, ``Sierra Popoluca''. \url{https://en.wikipedia.org/wiki/Sierra_Popoluca}
\end{itemize}

\section{Religious and Community Resources}
None found in search results.

\section{Digital and Audio Resources}
\begin{itemize}
\item Omniglot page on Sierra Popoluca (writing and orthography overview). \url{https://omniglot.com/writing/sierrapopoluca.htm}
\end{itemize}

\section{Curriculum Development Context}
The language is vigorous, but the search found no online dictionary or school materials. The 1999 dictionary and 1960 grammar are the main bases for orthography and vocabulary, and a more recent grammar (Boudreault 2018) is noted with a different orthography. Mexican bilingual education authorities and community teachers would be the natural partners, but no such programme was found in the results.

\end{document}